\begin{tikzpicture}[scale=1.0, >=stealth]
    \usetikzlibrary{patterns}
    % 定义坐标
    % 下方点 B (改为水平线)
    \coordinate (B1) at (0, 0);
    \coordinate (B2) at (1.5, 0);
    \coordinate (B3) at (3.0, 0);
    % 中间省略部分跨度 (距离减半，原为3.0，现为1.5)
    \coordinate (Bn) at (4.5, 0);
    \coordinate (Bn1) at (6.0, 0);

    % 上方点 A (A系列点搞紧凑一点，x坐标间隔减小，且往右上角微调)
    % A1 往右上角微调
    \coordinate (A1) at (1.0, 2.7);
    % A2
    \coordinate (A2) at (2.0, 2.9); 
    % A3
    \coordinate (A3) at (3.0, 3.1);
    
    % An 
    \coordinate (An) at (4.5, 3.4);
    % An+1 在 Bn+1 左上角 (Bn+1在6.0, An+1 x坐标要小于6.0)
    \coordinate (An1) at (5.5, 3.6);

    % 绘制两条主直线 (延伸出一点)
    \draw[thick] (-0.5, 0) -- (7.0, 0); % 下方直线 (水平)
    % 上方直线需要穿过 A 点
    % 斜率 k = (3.6 - 2.7) / (5.5 - 1.0) = 0.9 / 4.5 = 0.2
    % y - 2.7 = 0.2 * (x - 1.0) => y = 0.2x + 2.5
    % x = -0.5 => y = 2.4
    % x = 7.0 => y = 3.9
    \draw[thick] (-0.5, 2.4) -- (7.0, 3.9);   % 上方直线

    % 绘制三角形 S1
    \draw[thick] (A1) -- (B1);
    \draw[thick] (A1) -- (B2);
    \fill[pattern=north east lines, pattern color=black] (A1) -- (B1) -- (B2) -- cycle;
    \node at ($(A1)!0.7!($(B1)!0.5!(B2)$)$) {$S_1$};

    % 绘制三角形 S2
    \draw[thick] (A2) -- (B2);
    \draw[thick] (A2) -- (B3);
    \fill[pattern=north east lines, pattern color=black] (A2) -- (B2) -- (B3) -- cycle;
    \node at ($(A2)!0.7!($(B2)!0.5!(B3)$)$) {$S_2$};

    % 绘制 A3-B3 线段 (不形成完整三角形，因为后面是省略号)
    \draw[thick] (A3) -- (B3);

    % 绘制三角形 Sn
    \draw[thick] (An) -- (Bn);
    \draw[thick] (An) -- (Bn1);
    \fill[pattern=north east lines, pattern color=black] (An) -- (Bn) -- (Bn1) -- cycle;
    \node at ($(An)!0.7!($(Bn)!0.5!(Bn1)$)$) {$S_n$};

    % 绘制 An+1 - Bn+1 线段
    \draw[thick] (An1) -- (Bn1);

    % 标注点 A
    \node[above] at (A1) {$A_1$};
    \node[above] at (A2) {$A_2$};
    \node[above] at (A3) {$A_3$};
    \node[above] at (An) {$A_n$};
    \node[above] at (An1) {$A_{n+1}$};
    \node[above] at (6.7, 4) {$\cdots$}; % 上方直线的省略号
    \node at (3.8, 3.4) {$\cdots$}; % A3和An之间的省略号
    \node at (6.7, 1) {$\cdots$}; % 中间直线的省略号

    % 标注点 B
    \node[below] at (B1) {$B_1$};
    \node[below] at (B2) {$B_2$};
    \node[below] at (B3) {$B_3$}; 
    
    \node[below] at (Bn) {$B_n$};
    \node[below] at (Bn1) {$B_{n+1}$};
    \node at (6.7, -0.3) {$\cdots$}; % 下方直线的省略号
    \node at (3.75, -0.3) {$\cdots$}; % B3和Bn之间的省略号

    % 中间的省略号
    \node at (3.75, 1.0) {$\cdots$};

\end{tikzpicture}